\begin{appendices}

\section{Demonstração manual do dimensionamento}\label{ap:demonstracao}

Este apêndice reproduz, passo a passo e por cálculo manual, a verificação de uma solução da fronteira eficiente da célula de referência C-13 ($L = 5{,}0$~m, classe D40), de modo a permitir a conferência independente da implementação computacional.
Os cálculos foram conduzidos de forma inteiramente separada do código da plataforma, a partir apenas das equações apresentadas na Seção~\ref{sec:projeto_ponte}, e comparados ao final contra o resultado que a plataforma produz para a mesma solução.

\subsection{Solução verificada}

Escolheu-se a quarta das sete soluções listadas na Tabela~\ref{tab:solucoes_ref}, por ser uma solução intermediária da fronteira, nem o extremo de menor volume, nem o de menor flecha, com os valores de projeto na precisão plena com que foram gerados pela otimização:

\begin{align*}
d &= 67{,}094508~\text{cm}, & b_w &= 46{,}205444~\text{cm}, & h &= 10{,}174077~\text{cm}, \\
esp_{long} &= 108{,}475962~\text{cm}, & esp_{tab} &= 2{,}351631~\text{cm}. &&
\end{align*}

Os dois espaçamentos acima são as variáveis contínuas efetivamente manipuladas pelo NSGA-II. A rotina de acomodação geométrica descrita na Seção~\ref{sec:formulacao} converte-as no número inteiro de peças que cabem na pista e no vão, com o espaçamento uniforme resultante.
Para a longarina, distribuída ao longo da largura de pista $b_{w,\mathrm{pista}} = 4{,}5$~m, o número contínuo de peças é $n_{cont} = (b_{w,\mathrm{pista}} + esp_{long})/(d + esp_{long}) = 3{,}180922$, o que dá dois candidatos inteiros, $n = 3$ e $n = 4$.
Para $n = 3$, a folga $b_{w,\mathrm{pista}} - n\,d = 2{,}487165$~m reparte-se em dois vãos, resultando em $esp_{long,corr} = 1{,}243582$~m ($g_5 = -0{,}378209$). Para $n = 4$, a folga é $1{,}816220$~m, e o espaçamento corrigido cairia para $0{,}605407$~m ($g_5 = -0{,}697297$).
Ambos os candidatos são geometricamente admissíveis frente aos limites de $30$ a $200$~cm da Tabela~\ref{tab:intervalos}; prevalece então o critério de desempate pelo espaçamento mais próximo do valor contínuo original ($|1{,}243582 - 1{,}084760| = 0{,}158822$ contra $|0{,}605407 - 1{,}084760| = 0{,}479353$), e $n_{long} = 3$ é adotado, com $esp_{long,corr} = 124{,}358238$~cm.

Para o tabuleiro, distribuído ao longo do vão $L = 5{,}0$~m, o número contínuo é $n_{cont} = (L + esp_{tab})/(b_w + esp_{tab}) = 10{,}345591$, com candidatos $n = 10$ e $n = 11$.
Para $n = 11$, a folga $L - n\,b_w = 5{,}0 - 11 \times 0{,}462054 = -0{,}082599$~m já é negativa, isto é, as onze peças não cabem lado a lado no vão nem mesmo encostadas, e o candidato é descartado de imediato.
Resta apenas $n = 10$, com folga $0{,}379456$~m e espaçamento corrigido $esp_{tab,corr} = 4{,}216173$~cm, dentro dos limites de $2{,}0$ a $5{,}0$~cm da Tabela~\ref{tab:intervalos} ($g_6 = -0{,}156765$).
Ao contrário do que ocorre em outros pontos da fronteira, discutidos na Seção~\ref{sec:custo_robustez}, esta solução em particular não está próxima da restrição geométrica do tabuleiro.

\subsection{Propriedades geométricas}

Para a longarina circular, de diâmetro $d = 0{,}670945$~m, as Equações~\eqref{eq:areacirc}, \eqref{eq:icirc}, \eqref{eq:wcirc} e \eqref{eq:scirc} fornecem:
\begin{align*}
A_{long} &= \frac{\pi\,(0{,}670945)^2}{4} = 0{,}353561~\text{m}^2, \\
I_{long} &= \frac{\pi\,(0{,}670945)^4}{64} = 9{,}94759\times10^{-3}~\text{m}^4, \\
W_{long} &= \frac{I_{long}}{d/2} = 2{,}96525\times10^{-2}~\text{m}^3, \\
S_{long} &= \frac{(0{,}670945)^3}{12} = 2{,}51698\times10^{-2}~\text{m}^3.
\end{align*}

Para a prancha do tabuleiro, de largura $b_w = 0{,}462054$~m e altura $h = 0{,}101741$~m, as Equações~\eqref{eq:arearet} e~\eqref{eq:iret} fornecem:
\begin{align*}
A_{tab} &= 0{,}462054 \times 0{,}101741 = 4{,}70098\times10^{-2}~\text{m}^2, \\
I_{tab} &= \frac{0{,}462054 \times (0{,}101741)^3}{12} = 4{,}05506\times10^{-5}~\text{m}^4, \qquad W_{tab} = \frac{I_{tab}}{h/2} = 7{,}97135\times10^{-4}~\text{m}^3.
\end{align*}

\subsection{Resistências de cálculo}

Para a classe de carregamento de curta duração, madeira natural e classe de umidade 3 (Tabela~\ref{tab:parametros}), as Tabelas 4 e 5 da ABNT NBR 7190 fornecem $k_{mod,1} = 0{,}90$ e $k_{mod,2} = 0{,}80$, de modo que $k_{mod} = k_{mod,1}\,k_{mod,2} = 0{,}72$.
A partir dos valores característicos da classe D40 ($f_{m,k} = f_{c0,k} = 40{,}0$~MPa, conforme a equivalência do item 6.3.4 da \textcite{NBR7190-1:2022} discutida na Seção~\ref{sec:formulacao}, e $f_{v0,k} = 6{,}0$~MPa, Tabela~\ref{tab:classes_madeira}), a Equação~\eqref{eq:2.099} fornece as resistências de cálculo, iguais para a longarina e para o tabuleiro, que compartilham a mesma classe de resistência:
\begin{equation*}
f_{m,d} = \frac{40{,}0 \times 0{,}72}{1{,}40} = 20{,}571~\text{MPa}, \qquad f_{v0,d} = \frac{6{,}0 \times 0{,}72}{1{,}80} = 2{,}400~\text{MPa}.
\end{equation*}

\subsection{Ações}

O peso próprio da longarina segue a Equação~\eqref{eq:permanente}, com o peso específico da madeira $\gamma = 750 \times 9{,}81/1000 = 7{,}3575$~kN/m\textsuperscript{3}:
\begin{equation*}
G_{long} = \gamma\,A_{long} = 7{,}3575 \times 0{,}353561 = 2{,}60132~\text{kN/m}.
\end{equation*}

O peso próprio do tabuleiro é primeiro convertido em uma pressão equivalente, distribuindo o peso das $n_{tab} = 10$ peças sobre toda a área do tabuleiro ($b_{w,\mathrm{pista}} \times L$):
\begin{equation*}
p_{\mathrm{tab}} = \frac{\gamma\,n_{tab}\,(h\,b_w\,b_{w,\mathrm{pista}})}{b_{w,\mathrm{pista}}\,L} = \frac{7{,}3575 \times 10 \times (0{,}101741 \times 0{,}462054 \times 4{,}5)}{4{,}5 \times 5{,}0} = 0{,}691749~\text{kPa}.
\end{equation*}

Essa pressão, somada à carga permanente adicional $p_{gk} = 0{,}10$~kPa da Tabela~\ref{tab:parametros}, é distribuída sobre a faixa tributária de cada longarina, de largura $esp_{long,corr}$, e somada ao peso próprio da própria longarina:
\begin{equation*}
p_{gk,long} = (p_{gk} + p_{\mathrm{tab}})\,esp_{long,corr} + G_{long} = (0{,}10 + 0{,}691749)\times 1{,}243582 + 2{,}60132 = 3{,}58593~\text{kN/m}.
\end{equation*}

Para o tabuleiro, a mesma pressão combinada é convertida em carga linear ao longo do seu próprio vão local, distribuindo-a sobre a largura $b_w$ da peça:
\begin{equation*}
p_{gtab,k} = (p_{gk} + p_{\mathrm{tab}})\,b_w = (0{,}10 + 0{,}691749) \times 0{,}462054 = 0{,}365831~\text{kN/m}.
\end{equation*}

A carga variável de multidão, $p_{qk} = 4{,}0$~kPa da Tabela~\ref{tab:parametros}, é convertida em carga linear na longarina pela mesma faixa tributária, $p_{qk,long} = p_{qk}\,esp_{long,corr} = 4{,}0 \times 1{,}243582 = 4{,}97433$~kN/m.

Como o vão da longarina é $L = 5{,}0~\text{m} \leq 10{,}0$~m, a Equação~\eqref{eq:civ1} fornece $CIV = 1{,}35$, aplicado diretamente sobre os esforços variáveis, sem atenuação: o fator de $0{,}75$ que reduzia esse coeficiente em versões anteriores da plataforma não consta da \textcite{nbr7188} e foi removido (Seção~\ref{sec:formulacao}).
O mesmo coeficiente se aplica ao tabuleiro, cujo vão local (a distância entre longarinas, $esp_{long,corr} = 1{,}243582$~m) também é inferior a $10{,}0$~m, resultando no mesmo $CIV_{tab} = 1{,}35$.

\subsection{Esforços solicitantes e deslocamentos da longarina}

O momento fletor devido à carga permanente segue da Equação~\eqref{eq:mgk}:
\begin{equation*}
M_{gk} = \frac{3{,}58593 \times 5{,}0^2}{8} = 11{,}2060~\text{kN}\cdot\text{m}.
\end{equation*}

Como $L = 5{,}0$~m está no intervalo $3 < L \leq 6$~m, o momento devido à carga variável usa a Equação~\eqref{eq:mqk30}, sem a parcela de multidão, e é majorado diretamente pelo $CIV$:
\begin{equation*}
M_{qk,0} = \frac{3 \times 40 \times 5{,}0}{4} - 40 \times 1{,}5 = 90{,}000~\text{kN}\cdot\text{m}, \qquad M_{qk} = 90{,}000 \times 1{,}35 = 121{,}500~\text{kN}\cdot\text{m}.
\end{equation*}

O cortante devido à carga permanente segue da Equação~\eqref{eq:qgk}: $Q_{gk} = 3{,}58593 \times 5{,}0/2 = 8{,}96482$~kN.
Para o cortante variável (Equação~\eqref{eq:qqk}), a plataforma emprega, para seção circular, o próprio diâmetro da longarina no lugar do parâmetro geométrico genérico da fórmula, de modo que $e = L - 3a - 2d = 5{,}0 - 3\times1{,}5 - 2\times0{,}670945 = -0{,}841890$~m.
O valor negativo não invalida a expressão. Ele decorre apenas de a longarina ter diâmetro grande em relação ao vão e à distância entre eixos, e a fórmula é aplicada tal como definida, com $p_{qk,long} = 4{,}97433$~kN/m:
\begin{equation*}
Q_{qk,0} = \frac{40}{5{,}0}\big(6\times1{,}5 + 3\times(-0{,}841890)\big) + \frac{4{,}97433\times(-0{,}841890)^2}{2\times5{,}0} = 51{,}7946 + 0{,}3526 = 52{,}1472~\text{kN},
\end{equation*}
\begin{equation*}
Q_{qk} = 52{,}1472 \times 1{,}35 = 70{,}3987~\text{kN}.
\end{equation*}

As flechas seguem das Equações~\eqref{eq:flechag} e~\eqref{eq:flechaq}, com $E_{0,\mathrm{ef}} = 14{,}5$~GPa e $b = (L - 2a)/2 = 1{,}0$~m:
\begin{equation*}
\delta_{gk} = \frac{5 \times 3{,}58593 \times 5{,}0^4}{384 \times 14{,}5\times10^6 \times 9{,}94759\times10^{-3}} = 0{,}202~\text{mm},
\end{equation*}
\begin{equation*}
\delta_{qk} = \frac{40\,[5{,}0^3 + 2(1{,}0)(3\times5{,}0^2-4\times1{,}0^2)]}{48 \times 14{,}5\times10^6 \times 9{,}94759\times10^{-3}} = 1{,}543~\text{mm}.
\end{equation*}

\subsection{Verificações da longarina}

A combinação de ações no estado limite último (Equação~\eqref{eq:comb}) fornece o momento e o cortante de cálculo:
\begin{equation*}
M_{sd} = 11{,}2060 \times 1{,}30 + 121{,}500 \times 1{,}50 = 196{,}818~\text{kN}\cdot\text{m},
\end{equation*}
\begin{equation*}
Q_{sd} = 8{,}96482 \times 1{,}30 + 70{,}3987 \times 1{,}50 = 117{,}252~\text{kN}.
\end{equation*}

A verificação de flexão (Equação~\eqref{eq:verflexao}) resulta em:
\begin{equation*}
\sigma_{Mx,d} = \frac{196{,}818}{2{,}96525\times10^{-2}} = 6{,}637~\text{MPa},
\end{equation*}
\begin{equation*}
g_1 = \frac{6{,}637 - 20{,}571}{20{,}571} = -0{,}6773 \quad\text{(utilização de } 32{,}3\%\text{)}.
\end{equation*}

A verificação de cisalhamento, para seção circular, segue da Equação~\eqref{eq:vercis} com $S_{long} = 2{,}51698\times10^{-2}$~m\textsuperscript{3}:
\begin{equation*}
\tau_{\max} = \frac{Q_{sd}\,S_{long}}{I_{long}\,d} = \frac{117{,}252 \times 2{,}51698\times10^{-2}}{9{,}94759\times10^{-3} \times 0{,}670945} = 0{,}442~\text{MPa},
\end{equation*}
\begin{equation*}
g_2 = \frac{0{,}442 - 2{,}400}{2{,}400} = -0{,}8158 \quad\text{(utilização de } 18{,}4\%\text{)}.
\end{equation*}

Para a flecha, a Equação~\eqref{eq:verflecha} compara duas combinações em paralelo, cada qual com seu próprio limite normativo (Tabela 21 da \textcite{NBR7190-1:2022}).
Na combinação quase permanente, a fluência amplia as duas parcelas, permanente e variável (esta última já reduzida por $\psi_2$): $\delta_{tot} = (1+\varphi)\,(\delta_{gk} + \psi_2\,\delta_{qk}) = 1{,}80 \times (0{,}202 + 0{,}30 \times 1{,}543) = 1{,}80 \times 0{,}665 = 1{,}197$~mm.
Contra o limite $\delta_{lim,total} = L/350 = 14{,}286$~mm, a utilização é de $8{,}4\%$ ($g_{3a} = -0{,}9162$).
Já na combinação rara, exigida em paralelo, verifica-se a flecha instantânea da parcela variável isolada, sem fluência, contra um limite mais severo: $\delta_{qk} = 1{,}543$~mm contra $\delta_{lim,var} = L/500 = 10{,}000$~mm, utilização de $15{,}4\%$ ($g_{3b} = -0{,}8457$), que é a mais severa das duas e portanto governa: $g_3 = \max(g_{3a}, g_{3b}) = -0{,}8457$.
O fator de flecha $f_2 = \delta_{tot}/(L/350)$ reportado nas demais tabelas do artigo refere-se sempre à combinação quase permanente, $f_2 = 0{,}0838$.

\subsection{Verificação do tabuleiro}

A prancha do tabuleiro vence, como viga biapoiada, o vão local $esp_{long,corr} = 1{,}243582$~m entre longarinas.
O momento devido à carga permanente segue da Equação~\eqref{eq:mgk}: $M_{gk,tab} = 0{,}365831 \times 1{,}243582^2/8 = 0{,}0707$~kN$\cdot$m.
O momento devido à carga variável decorre da roda do trem tipo aplicada sobre esse vão local, com comprimento efetivo de contato $a_r = 0{,}45$~m, e é majorado diretamente pelo $CIV_{tab}$:
\begin{equation*}
M_{qk,tab,0} = \frac{40}{4}\,(1{,}243582 - 0{,}45) = 7{,}9358~\text{kN}\cdot\text{m},
\end{equation*}
\begin{equation*}
M_{qk,tab} = 7{,}9358 \times 1{,}35 = 10{,}7134~\text{kN}\cdot\text{m}.
\end{equation*}

O momento de cálculo e a verificação de flexão resultam em:
\begin{equation*}
M_{sd,tab} = 0{,}0707\times1{,}30 + 10{,}7134\times1{,}50 = 16{,}1620~\text{kN}\cdot\text{m},
\end{equation*}
\begin{equation*}
\sigma_{Mx,d,tab} = \frac{16{,}1620}{7{,}97135\times10^{-4}} = 20{,}275~\text{MPa},
\end{equation*}
\begin{equation*}
g_4 = \frac{20{,}275 - 20{,}571}{20{,}571} = -0{,}0144 \quad\text{(utilização de } 98{,}6\%\text{)}.
\end{equation*}

\subsection{Volume total}

Pela Equação~\eqref{eq:f1}, com $n_{long} = 3$ e $n_{tab} = 10$:
\begin{equation*}
V_{\mathrm{total}} = 3\times0{,}353561\times5{,}0 + 10\times4{,}70098\times10^{-2}\times4{,}5 = 5{,}30341 + 2{,}11544 = 7{,}41885~\text{m}^3.
\end{equation*}

\subsection{Síntese comparativa}\label{ap:sintese}

A Tabela~\ref{tab:ap_sintese} confronta os oito resultados obtidos manualmente com os retornados pela plataforma para a mesma solução, no ponto nominal, isto é, sem a média sobre as $N_c = 30$ realizações perturbadas da avaliação robusta (Seção~\ref{sec:robustez}).
Essa distinção é deliberada: a robustez é um procedimento numérico de agregação estatística sobre repetições do mesmo modelo mecânico, e não faz parte do que este apêndice verifica.
Por isso, os valores de referência aqui não coincidem com o $g_{\max}$ publicado na Tabela~\ref{tab:solucoes_ref} para a mesma solução, que é a média robusta.

\begin{table}[!ht]
\caption{Comparação entre o cálculo manual e a avaliação nominal da plataforma, para a solução da Seção~\ref{ap:demonstracao}.\label{tab:ap_sintese}}
\begin{threeparttable}
\footnotesize
\setlength{\tabcolsep}{3pt}
\begin{tabular*}{\columnwidth}{@{\extracolsep\fill}lccc@{\extracolsep\fill}}
\toprule
Grandeza & Manual & Plataforma\tnote{a} & Diferença \\
\midrule
$V_{\mathrm{total}}$ (m\textsuperscript{3}) & 7{,}41885 & 7{,}41885 & $< 0{,}001\%$ \\
$f_2 = \delta_{tot}/(L/350)$ & 0{,}0838 & 0{,}0838 & $< 0{,}001\%$ \\
$g_1$: flexão da longarina & $-0{,}6773$ & $-0{,}6773$ & $< 0{,}001\%$ \\
$g_2$: cisalhamento da longarina & $-0{,}8158$ & $-0{,}8158$ & $< 0{,}001\%$ \\
$g_3$: flecha da longarina & $-0{,}8457$ & $-0{,}8457$ & $< 0{,}001\%$ \\
$g_4$: flexão do tabuleiro & $-0{,}0144$ & $-0{,}0144$ & $< 0{,}001\%$ \\
$g_5$: espaçamento entre longarinas & $-0{,}3782$ & $-0{,}3782$ & $< 0{,}001\%$ \\
$g_6$: espaçamento entre peças do tabuleiro & $-0{,}1568$ & $-0{,}1568$ & $< 0{,}001\%$ \\
\bottomrule
\end{tabular*}
\begin{tablenotes}
\footnotesize
\item[a] Avaliação nominal, sem a média sobre as realizações perturbadas da avaliação robusta (ver texto).
\item As diferenças remanescentes, na quinta casa decimal ou além, são atribuíveis exclusivamente ao arredondamento das variáveis de entrada nesta demonstração e não a divergência entre o memorial e a implementação.
\end{tablenotes}
\end{threeparttable}
\end{table}

A conferência confirma que o modelo mecânico documentado na Seção~\ref{sec:projeto_ponte} reproduz fielmente a implementação computacional para as oito grandezas verificadas.
Cabe registrar a natureza da pequena diferença entre o $g_{\max}$ nominal desta demonstração e o valor publicado na Tabela~\ref{tab:solucoes_ref} para a mesma linha: aqui, o maior valor entre as seis funções limite é $g_4 = -0{,}0144$ (flexão do tabuleiro, ponto nominal, sem perturbação), ao passo que a tabela publica $g_{\max} = -0{,}0097$, a média das seis funções limite sobre as $N_c = 30$ realizações perturbadas da avaliação robusta (Seção~\ref{sec:robustez}).
As duas leituras concordam quanto à verificação governante, a flexão do tabuleiro, mas não precisam coincidir em valor: a primeira descreve o comportamento da solução exatamente no ponto de projeto, a segunda descreve seu comportamento esperado sob a variabilidade dimensional considerada na busca.
Esse é o mesmo mecanismo quantificado em escala na Seção~\ref{sec:custo_robustez}: a robustez desloca o critério de aceitação do ponto nominal para o comportamento médio sob incerteza, de modo que as duas avaliações --- nominal e robusta --- são complementares, e não intercambiáveis, na leitura de uma solução da fronteira.

\end{appendices}
